% ----------------------
\subsection{Section 95 (1 June 1833)}
% ----------------------
	\label{DC95}
	
	% -----
	\subsubsection{Historical Background}
	% -----
		\label{DC95-HB}
		
		% --
		\paragraph{JSP}
		% --
		
			``In late December 1832 and early January 1833, as part of a call to educate men for the ministry, a revelation instructed church members to organize a `school of the prophets' and to build a House of the Lord wherein individuals would learn the law and receive power that had been previously promised. The revelation further told them that they should `prepare evry needful [thing]' for the school and that the house should be `an house of prayer, an house of fasting, an house of faith, an house of Learning, an house of glory, an house of order an house of God.' During the early months of 1833, before construction on the house began, the first School of the Prophets, which involved just over a dozen church elders, met in Kirtland, Ohio. On 11 January 1833, JS emphasized the urgency of building the house and holding the school in a letter to William W. Phelps in Jackson County, Missouri. JS stated, `You will see that the Lord commanded us in Kirtland to build an house of God, \& establish a school for the Prophets, this is the word of the Lord to us, \& we must--- yea the Lord helping us we will obey, as on conditions of our obedience, he has promised us great things, yea even a visit from the heavens to honor us with his own presence.'
			
			``On 4 May, nearly four months after JS wrote about the urgent need to build the house, a committee was formed to raise funds for the building's construction. A conference unanimously voted that a committee, consisting of Hyrum Smith, Jared Carter, and Reynolds Cahoon, be formed to raise the money. Nearly a month later, neither this committee nor other church leaders had taken any steps toward constructing the building.
			
			``The revelation featured here, dictated on 1 June, stated that church leaders' unresponsiveness to the call to build this religious structure would halt the growth and spiritual work of the church. The revelation also specified the building's dimensions and directed the manner in which church members should construct it. Immediately following this revelation, a conference of high priests discussed the revelation and began drafting construction plans. The conference appointed JS, Sidney Rigdon, and Frederick G. Williams---the presidency of the high priesthood---to serve as the planning committee, which was responsible for obtaining `a draft or construction of the inner court of the house.'
			
			``Likely in response to this revelation, the building committee, which still comprised Hyrum Smith, Jared Carter, and Reynolds Cahoon, prepared a circular for the members of the Church of Christ that same day, writing that they had been officially charged with soliciting subscriptions to establish a fund to build the house. `Unless we fulfil this command vis establish an house,' they warned, `and prepare all things necessary whereby the Elders may gather into a school called the school of the prophets and receive that instruction that the Lord designs they should receive we may all dispare of obtaining the great blessing that God has promised to the faithful of the Church of Christ.' The circular reiterated the promise found in the 2 January 1831 revelation that God would endow individuals with `power from on high,' which they would gain from `that instruction that the Lord designs' in the School of the Prophets. The committee encouraged church members to `make evry possable exertion to aid temporally as well as spiritually in this great work that the Lord is bringing about and is about to accomplish.' The circular also called for church members to pay their subscriptions and to send the funds to Kirtland by 1 September 1833. Hyrum Smith, Reynolds Cahoon, and others began construction on the multipurpose House of the Lord by 7 June 1833. The House of the Lord, completed in 1836, eventually served as both a school and a place of worship; it was the church's first temple.''
			
	% -----
	\subsubsection{Versions}
	% -----
		\label{DC95-V}
		
		See the \href{https://drive.google.com/file/d/1goAvNz8jS4R8slYfMG_db55Ty2z_vmgK/view?usp=sharing}{sections comparison} document.
	
		\begin{compactdesc}
			\item[JSP] \hyperref[EMS-1833-JS-1-JUN-1833]{1 June 1833}
			\item[BC] ---
			\item[DC 1835] Section 95, 233--234
			\item[DC 1844] Section 96, 362--364
			\item[TS] 6:2 (1 February 1845), 784
			\item[MS] 14:28 (4 September 1852), 436
		\end{compactdesc}

	% -----
	\subsubsection{Section 95:1} % *DC95-1 
	% -----
		\label{DC95-1}

		VERILY, thus saith the Lord unto you whom I love, and whom I love I also chasten that their sins may be forgiven, for with the chastisement I prepare a way for their deliverance in all things out of temptation, and I have loved you---

		% \begin{framed}
		% \scriptsize
			% \begin{compactdesc}
				% \item[JSP] 
				% % \item[BC] 
				% % \item[]
			% \end{compactdesc}
		% \end{framed}

	% -----
	\subsubsection{Section 95:2} % *DC95-2 
	% -----
		\label{DC95-2}

		Wherefore, ye must needs be chastened and stand rebuked before my face;

		% \begin{framed}
		% \scriptsize
			% \begin{compactdesc}
				% \item[JSP] 
				% % \item[BC] 
				% % \item[]
			% \end{compactdesc}
		% \end{framed}

	% -----
	\subsubsection{Section 95:3} % *DC95-3 
	% -----
		\label{DC95-3}

		For ye have sinned against me a very grievous sin, in that ye have not considered the great commandment in all things, that I have given unto you concerning the building of mine house;

		% \begin{framed}
		% \scriptsize
			% \begin{compactdesc}
				% \item[JSP] 
				% % \item[BC] 
				% % \item[]
			% \end{compactdesc}
		% \end{framed}

	% -----
	\subsubsection{Section 95:4} % *DC95-4 
	% -----
		\label{DC95-4}

		For the preparation wherewith I design to prepare mine apostles to prune my vineyard for the last time, that I may bring to pass my strange act,%
			\footnote{%
				% \label{}%
				See \hyperref[ISAIAH-28-21]{Isaiah 28:21}.
				}
		that I may pour out my Spirit upon all flesh---

		% \begin{framed}
		% \scriptsize
			% \begin{compactdesc}
				% \item[JSP] 
				% % \item[BC] 
				% % \item[]
			% \end{compactdesc}
		% \end{framed}

	% -----
	\subsubsection{Section 95:5} % *DC95-5 
	% -----
		\label{DC95-5}

		But behold, verily I say unto you, that there are many who have been ordained among you, whom I have called but few of them are chosen.

		% \begin{framed}
		% \scriptsize
			% \begin{compactdesc}
				% \item[JSP] 
				% % \item[BC] 
				% % \item[]
			% \end{compactdesc}
		% \end{framed}

	% -----
	\subsubsection{Section 95:6} % *DC95-6 
	% -----
		\label{DC95-6}

		They who are not chosen have sinned a very grievous sin, in that they are walking in darkness at noon-day.

		% \begin{framed}
		% \scriptsize
			% \begin{compactdesc}
				% \item[JSP] 
				% % \item[BC] 
				% % \item[]
			% \end{compactdesc}
		% \end{framed}

	% -----
	\subsubsection{Section 95:7} % *DC95-7 
	% -----
		\label{DC95-7}

		And for this cause I gave unto you a commandment that you should call your solemn assembly, that your fastings and your mourning might come up into the ears of the Lord of \hyperref[SABAOTH]{Sabaoth}, which is by interpretation, the creator of the first day, the beginning and the end.

		\begin{framed}
		\scriptsize
			\begin{compactdesc}
				\item[JSP] and for this cause I gave unto you a commandment that you should call your solem assembly that your fastings and your mourning might come up into the ears of the Lord of sabaoth which is by interpretation the creator of \sout{all things} the first day the beginning and the end.
				% \item[BC] 
				% \item[]
			\end{compactdesc}
		\end{framed}

	% -----
	\subsubsection{Section 95:8} % *DC95-8 
	% -----
		\label{DC95-8}

		Yea, verily I say unto you, I gave unto you a commandment that you should build a house, in the which house I design to endow those whom I have chosen with power from on high;

		% \begin{framed}
		% \scriptsize
			% \begin{compactdesc}
				% \item[JSP] 
				% % \item[BC] 
				% % \item[]
			% \end{compactdesc}
		% \end{framed}

	% -----
	\subsubsection{Section 95:9} % *DC95-9 
	% -----
		\label{DC95-9}

		For this is the promise of the Father unto you; therefore I command you to tarry, even as mine apostles at Jerusalem.

		% \begin{framed}
		% \scriptsize
			% \begin{compactdesc}
				% \item[JSP] 
				% % \item[BC] 
				% % \item[]
			% \end{compactdesc}
		% \end{framed}

	% -----
	\subsubsection{Section 95:10} % *DC95-10 
	% -----
		\label{DC95-10}

		Nevertheless, my servants sinned a very grievous sin; and contentions arose in the school of the prophets; which was very grievous unto me, saith your Lord; therefore I sent them forth to be chastened.

		% \begin{framed}
		% \scriptsize
			% \begin{compactdesc}
				% \item[JSP] 
				% % \item[BC] 
				% % \item[]
			% \end{compactdesc}
		% \end{framed}

	% -----
	\subsubsection{Section 95:11} % *DC95-11 
	% -----
		\label{DC95-11}

		Verily I say unto you, it is my will that you should build a house.  If you keep my commandments you shall have power to build it.

		% \begin{framed}
		% \scriptsize
			% \begin{compactdesc}
				% \item[JSP] 
				% % \item[BC] 
				% % \item[]
			% \end{compactdesc}
		% \end{framed}

	% -----
	\subsubsection{Section 95:12} % *DC95-12 
	% -----
		\label{DC95-12}

		If you keep not my commandments, the love%
			\footnote{%
				% \label{}%
				The note at the end of this verse understands this ``love'' as the love the Father has \textit{for us}---that doesn't continue with us, and as a result, we are walking in darkness. That's a subjective-genitive interpretation of this construction, but it's also possible to go objective-genitive, where the love we have \textit{for the Father} no longer continues with us, and as a result, we are walking in darkness. See also \hyperref[JOHN-15-9]{John 15:9--10}, \hyperref[1JOHN-2-15]{1 John 2:15} (this one could also go either way, but by parallel with ``love the world,'' it could be objective genitive; see notes there, ZECNT says basically the same thing), \hyperref[MOSIAH-2-4]{Mosiah 2:4}, 
				}
		of the Father shall not continue with you, therefore you shall walk in darkness.%
			\footnote{%
				% \label{}%
				Very interesting verse; see \hyperref[DC95-6]{verse 6} above. See also \hyperref[2NEPHI-4-21]{2 Nephi 4:21} and notes there, as this scripture is a key in understanding God's light/power/love/spirit as all the same thing, or those are words for referring to the same thing but calling attention to different aspects of its activity.
				}

		% \begin{framed}
		% \scriptsize
			% \begin{compactdesc}
				% \item[JSP] 
				% % \item[BC] 
				% % \item[]
			% \end{compactdesc}
		% \end{framed}

	% -----
	\subsubsection{Section 95:13} % *DC95-13 
	% -----
		\label{DC95-13}

		Now here is wisdom, and the mind of the Lord---let the house be built, not after the manner of the world, for I give not unto you that ye shall live after the manner of the world;

		% \begin{framed}
		% \scriptsize
			% \begin{compactdesc}
				% \item[JSP] 
				% % \item[BC] 
				% % \item[]
			% \end{compactdesc}
		% \end{framed}

	% -----
	\subsubsection{Section 95:14} % *DC95-14 
	% -----
		\label{DC95-14}

		Therefore, let it be built after the manner which I shall show unto three of you, whom ye shall appoint and ordain unto this power.

		% \begin{framed}
		% \scriptsize
			% \begin{compactdesc}
				% \item[JSP] 
				% % \item[BC] 
				% % \item[]
			% \end{compactdesc}
		% \end{framed}

	% -----
	\subsubsection{Section 95:15} % *DC95-15 
	% -----
		\label{DC95-15}

		And the size thereof shall be fifty and five feet in width, and let it be sixty-five feet in length, in the inner court thereof.

		% \begin{framed}
		% \scriptsize
			% \begin{compactdesc}
				% \item[JSP] 
				% % \item[BC] 
				% % \item[]
			% \end{compactdesc}
		% \end{framed}

	% -----
	\subsubsection{Section 95:16} % *DC95-16 
	% -----
		\label{DC95-16}

		And let the lower part of the inner court be dedicated unto me for your sacrament offering, and for your preaching, and your fasting, and your praying, and the offering up of your most holy desires unto me, saith your Lord.

		% \begin{framed}
		% \scriptsize
			% \begin{compactdesc}
				% \item[JSP] 
				% % \item[BC] 
				% % \item[]
			% \end{compactdesc}
		% \end{framed}

	% -----
	\subsubsection{Section 95:17} % *DC95-17 
	% -----
		\label{DC95-17}

		And let the higher part of the inner court be dedicated unto me for the school of mine apostles, saith Son Ahman;%
			\footnote{%
				% \label{}%
				See the ``\hyperref[EMS-1832-JS-CIR-4-MAR-1832-LANG]{Sample of pure language}'' from early 1832. 
				}
		or, in other words, Alphus;%
			\footnote{%
				% \label{}%
				See \hyperref[REVELATIONNT-1-8]{Revelation 1:8} and notes there.
				}
		or, in other words, Omegus;%
			\footnote{%
				% \label{}%
				This string of titles is very bizarre---how have Latin noun endings been attached to these Greek letter names? Makes very little sense, but the JSP version helps answer the question, I think (note that the JSP is RB2). The JSP/RB2 and RB1 accounts have ``Alphas/Omegas'' and ``Alpha/Omegas,'' respectively; the ``Alphus/Omegus'' stuff somehow gets introduced in the 1835 version of the DC. I don't know why this would have happened---``Alphas'' and ``Omegas'' just could have been plural, and that seems to be the correct way to understand what is really happening in this verse with these titles for Christ.
				}
		even Jesus Christ your Lord. Amen.

		\begin{framed}
		\scriptsize
			\begin{compactdesc}
				\item[JSP] and let the higher part of the inner court be dedicated unto me for the school of mine Apostles saith Son Ah Man, or in otherwords Alphas, or in other words Omegas even Jesus Christ your lord Amen.------
				\item[RB1] And let the higher part of the inner court be dedicated unto me for the school of mine apostles, saith Son awman, or in other words, Alpha, or in other words Omegas, even Jesus Christ your Lord: Amen.
				\item[DC 1835] And let the higher part of the inner court, be dedicated unto me for the school of mine apostles, saith Son Ahman; or, in other words, Alphus; or, in other words, Omegus; even Jesus Christ your Lord. Amen.
				\item[DC 1844] And let the higher part of the inner court, be dedicated unto me for the school of mine apostles, saith Son Ahman; or, in other words, Alphus; or, in other words, Omegus; even Jesus Christ your Lord: Amen.
			\end{compactdesc}
		\end{framed}